% SPDX-License-Identifier: MIT
\chapter{Terminal, Rendered und Portabilitaet}\label{chap:backends}
\ChapterSummary{Dieses Kapitel hilft bei der Auswahl eines
Ausgabebackends. Sie lernen die praktischen Unterschiede
zwischen Terminalzellen, Grafikpixeln und nativen Widgets
kennen sowie die Sicherheits- und Testfragen, die gerade
fuer Administrationssoftware entscheidend sind.}

\section{Terminal als vollwertiges Ziel}
Die Terminalfamilie ist nicht lediglich ein
schlechter Ersatz fuer eine Desktop-GUI.
Sie bietet Bedienung ueber SSH, geringe
grafische Abhaengigkeiten, klare Tastaturwege
und Moeglichkeiten fuer Systeme ohne
Desktop-Umgebung. Genau das ist fuer
Remoteadministration und Wartung
interessant.

Die aktuelle Terminal-Schicht enthaelt
unter anderem \term{ScreenBuffer},
\term{TextMetrics}, \term{AnsiFrameEncoder},
\term{TerminalSession},
\term{TerminalDevice},
\term{TerminalBackend} und
\term{TerminalPresentationSink}.
Die normale Widget-Semantik bleibt im Core.

\section{Der Weg von Eingabe zu Bild}
\begin{enumerate}
\item Ein Terminal liefert Bytes oder eine
plattformabhaengige Konsoleneingabe.
\item Der Backend-Adapter ordnet daraus
Tasten-, Text-, Resize- und gegebenenfalls
Pointer-Ereignisse zu.
\item Fokus, Eventrouting, Command oder Widget
aendern den semantischen Zustand.
\item Messung und Praesentation erzeugen
einen Bildschirmzustand im \term{ScreenBuffer}.
\item Der \term{AnsiFrameEncoder} wandelt
den fertigen Frame in ANSI/VT-Bytes um;
das Device sendet sie an das Terminal.
\end{enumerate}
Die Trennung ist absichtlich:
Ein Widget darf aus beliebigem Anwendungstext
nicht direkt rohe Terminal-Control-Sequenzen
an das Geraet senden. Sicherheitskritischer
Text muss Dateninhalt bleiben.

\section{Zeichenzellen sind keine Bytes}
Ein UTF-8-String kann aus vielen Bytes,
weniger Unicode-Scalars und einer ganz
anderen Anzahl Terminalzellen bestehen.
Besonders breite ostasiatische Zeichen,
kombinierende Zeichen und unterschiedliche
Terminal-Font-/Breitenregeln sind relevant.
\term{TextMetrics} besitzt dafuer eine
explizite Breitenpolitik und
Wide-Cell-Belegung.
Eine Anwendung sollte niemals mit
\texttt{string.size()} die sichtbare
Terminalbreite berechnen.

\section{RAII-Sitzung und Wiederherstellung}
Waehrend einer Terminalanwendung koennen
Raw-Modus, alternativer Bildschirm,
VT-Modi und gegebenenfalls xterm-
Pointer-Reporting aktiviert werden.
Die \term{TerminalSession} kapselt
Erwerb und Freigabe dieser Geraeteressourcen.
Das reduziert die Gefahr, nach regularem
Programmende eine unbrauchbare Konsole
zurueckzulassen.

\Important{Wenn Sie eine Anwendung mit einem
harten Kill abbrechen, etwa bei einem
Systemabsturz oder einem nicht abfangbaren
Signal, kann keine C++-Destruktion garantiert
werden. Testen Sie deshalb auch den
Wiederherstellungspfad Ihres konkreten Hosts,
und verwenden Sie sichere Fallbacks fuer
Notfallsituationen.}

\section{Pointer-Protokolle und Tastatur-Fallback}
Moderne Terminals koennen Mauspositionen,
Klicks, Bewegung und Scrollsignale ueber
xterm-kompatible Protokolle liefern.
Die Terminal-Demo kann
\term{TerminalPointerTrackingMode}
bewusst setzen. Alte VT-artige Terminals
muessen das jedoch nicht unterstuetzen.
Kernaktionen sollen deshalb ueber
Fokus, Enter/Space, Escape oder
erreichbare Commands moeglich bleiben.
Die reine Existenz eines PointerHandlers
ist keine portable Bediengarantie.

\section{Rendered: Befehlsliste statt SDL in Core}
\term{DisplayList} und \term{RenderDevice}
ermoeglichen es, Grafikoperationen
deterministisch zu beschreiben und
zu wiederholen. Die
\term{RenderedMeasurementContext}
liefert dazu Schriftmetriken.
\term{RenderedPresentationSink}
uebersetzt semantische Widget-Zustaende
in Zeichenvorgaenge. Diese API ist
nicht gleich SDL3 und kann ohne
graphische Sitzung getestet werden.

Erst der optionale \term{SDL3}-Adapter
realisiert eine konkrete Fenster- oder
Headless-Grafikumgebung mit SDL3/SDL\_ttf.
Er ist experimentell, opt-in und
noch kein garantiert installierbarer
externer Adapter fuer beliebige Projekte.
Ein Anwendungsautor sollte einen
Backendwechsel nicht als
ungetesteten Umschalter behandeln.

\section{Native Peers sind spaeter vorgesehen}
Eine native Win32-, GTK- oder AppKit-
Implementierung waere eine weitere
Darstellung derselben Widgets,
nicht eine neue Semantik. Es ist aber
\emph{kein} fertiger Paritaetsstand
vorhanden. Plattformintegration,
Barrierefreiheit, Text-Input-
Methoden, HiDPI, Fensterlebensdauer
und native Menues sind eigene
Entwicklungsaufgaben.

\section{Testmatrix fuer portable Programme}
\begin{tabularx}{\linewidth}{@{}p{35mm}X@{}}
\toprule
Testbereich & Zu pruefendes Verhalten\\\midrule
Terminal klein/gross & Umbruch, Clipping,
Fokus, Popup-Position, Resize.\\
UTF-8 & Mehrbyte-Zeichen, breite
Zellen, Auswahl, Cursor, Loeschen.\\
Tastatur & Tab, Shift+Tab, Escape,
Funktionstasten, Modifikatoren.\\
Pointer & Capture, Fenster verlassen,
Popup-Ziel, Doppelklick und Ziehen.\\
Lebensdauer & Widget entfernt, Model
reset, Subscription getrennt, Callback-
Zerstoerung.\\
Backend & Headless-Core-Tests,
Terminal-PTY/ConPTY, bei Bedarf SDL3.\\
\bottomrule
\end{tabularx}
Die vorhandene CI und die manuellen
Smoke-Test-Protokolle unter
\filepath{docs/de/} und \filepath{docs/en/}
sind dafuer Referenzen.

\PracticalTip{Fuer ein Administrationsprogramm
ist ein vollstaendiger, tastaturbedienbarer
Terminalpfad oft wertvoller als eine
visuell beeindruckende, aber noch
unvollstaendige SDL3-Version.}
